\answer{ex:21:1}{(a)~$T=\tau_{\text{rec}}\ln\frac{1}{1-x^*}$: $1.84$~s와 $3.68$~s. (b)~$3.36$에서 $4.24$~s까지. 민감도 $\dd T/\dd x^*=\tau_{\text{rec}}/(1-x^*)=80$~s이므로 $x^*\to1$에서 일제 발화는 불규칙해진다.}
\answer{ex:21:2}{(a)~\eqref{eq:energy}에서처럼 $8.6$~W. (b)~$205$~Mbit/s, $0.2$~W로 배양($10^{-4}$~W)의 $2\cdot10^3$배이다.}
\answer{ex:21:3}{$x_{\text{ss}}=1/(1+Ur_b\tau_{\text{rec}})<x^*$이려면 $r_b>(1/x^*-1)/(U\tau_{\text{rec}})$, 즉 $x^*=0.9$에서 $0.28$~Hz, $x^*=0.99$에서 $0.025$~Hz. 시냅스 세기는 $x_{\text{ss}}$로, 즉 $10\%$와 $1\%$만큼 떨어진다.}
\answer{ex:21:4}{$u(t-k)$는 정규직교이고 $m_k=\|P_Vu(t-k)\|^2$이다. 여기서 $P_V$는 출력이 생성하는 $N$차원 공간으로의 사영이다. 베셀 부등식에 의해 $\sum_k m_k\le N$.}
\answer{ex:21:5}{$\tanh$가 있으면 $\mathrm{MC}\approx9$--$11$, 선형 저수지는 $15$--$18$(시드 세 개), 둘 다 $\le30$이다. 비선형성의 대가는 기억이다.}
\answer{ex:21:6}{(a)~$n\le102$. (b)~$5.6$시그마.}
\answer{ex:21:7}{열린 문제. 핵심 대조 실험: 판독을 고정하고, 자극 없는 휴지 뒤에도 향상이 남는지 확인한다. 남지 않으면 바뀐 것은 가중치가 아니라 상태이다.}
